\chapter{Cotangent deformations and their generating bundles}
\label{chap:cs-surfaces}

The quadratic cover \(z^2=t\) explains the two modules at fusion.
It does not yet provide a surface on which both modules occur.
For that purpose we add a coordinate \(p\), interpreted as the linear
coordinate in a cotangent direction. The involution must change the
sign of both coordinates: this makes \(dz\wedge dp\) invariant.
The resulting algebra admits a deformation
\[
 pz-zp=\lambda g,\qquad g^2=1.
\]
Its centre will be the surface \(XY=W^2-\lambda^2/4\).
For nonzero \(\lambda\), we shall construct an equivalence with modules
on that surface, including the vector bundle on which the original
Yangian acts. At zero the equivalence fails for a specific reason:
one of the character modules is killed by the proposed generator.

The same calculation has a cyclic version. It will show exactly when
a chosen character generates all modules, and will exhibit a missing
module whenever it does not. This turns the smoothness condition into
a representation-theoretic condition with an explicit converse.

\section{The cotangent algebra and its zero section}

Work first over \(\mathbb C\). Let
\[
 A=\mathbb C[z]\rtimes\langle g\rangle,\qquad
 B_0=\mathbb C[z,p]\rtimes\langle g\rangle,
 \qquad g^2=1,\quad gz=-zg,\quad gp=-pg.
\]
The notation for a skew group algebra was defined in
Chapter~\ref{chap:fc-fusion}. In \(B_0\), \(z,p\) commute.
The inclusion \(A\to B_0\) and quotient \(B_0\to A\) sending \(p\)
to zero are actual algebra maps. Geometrically the latter is restriction
to the zero section. It does not imply that derived endomorphisms of
that section are concentrated in degree zero.

Put \(e_\pm=(1\pm g)/2\) and define arrows
\[
 a=e_+ze_-,\quad b=e_-ze_+,\qquad
 c=e_+pe_-,\quad d=e_-pe_+.
\]
Then \(z=a+b\), \(p=c+d\), and the commutation relation is precisely
\begin{equation}\label{eq:cs-preprojective}
 cb-ad=0,\qquad da-bc=0.
\end{equation}
Conversely these equations and the endpoint rules reconstruct \(B_0\).
They are the preprojective relations for the graph with two vertices
and two edges; the four arrows and their signs specify the convention.
Here a path algebra has as basis all finite composable words in its
arrows, with vertices as constant paths. Relations mean the two-sided
ideal generated by the displayed differences of paths.

Let \(P=B_0/(p)\), considered as a right \(B_0\)-module. The notation
\((p)\) denotes the two-sided ideal, which is also the right ideal
\(pB_0\), because \(p\) commutes with \(z\) and anticommutes with \(g\).
A projective resolution is
\begin{equation}\label{eq:cs-section-resolution}
 \mathsf P=(B_0\xrightarrow{L_p}B_0),
\end{equation}
in degrees \(-1,0\). Left multiplication by \(p\) is right-linear
and injective on the polynomial basis.

Define an automorphism \(\nu\) of \(A\) by \(\nu(z)=z\),
\(\nu(g)=-g\). The equation \(ap=p\nu(a)\) holds for every \(a\in A\).
It gives a chain endomorphism \(F_a\) of \(\mathsf P\), acting by
\(L_{\nu(a)}\) in degree \(-1\) and by \(L_a\) in degree zero.
Let \(\epsilon\) be the degree-one endomorphism which is the identity
from the degree-\(-1\) copy of \(B_0\) to the degree-zero copy.

\begin{proposition}\label{prop:cs-normal-extension}
The derived endomorphism algebra of \(P\) has the differential graded
model
\[
 E=A\langle\epsilon\rangle/
       (\epsilon^2,\ \epsilon a-\nu(a)\epsilon),\qquad
 |\epsilon|=1,\qquad d=0.
\]
In particular, the normal self-extension is nonzero and carries the
sign character.
\end{proposition}
\begin{proof}
The maps \(F_a\) multiply as \(F_aF_b=F_{ab}\) and commute with
the resolution differential. Direct composition gives
\(\epsilon F_a=F_{\nu(a)}\epsilon\) and \(\epsilon^2=0\).
They therefore define a differential graded algebra map
\(E\to\operatorname{End}_{B_0}(\mathsf P)\).
Postcomposition with \(\mathsf P\to P\) gives a quasi-isomorphism of
complexes to \(\operatorname{Hom}_{B_0}(\mathsf P,P)\).
To justify this last assertion, its cone is
\(\operatorname{Hom}_{B_0}(\mathsf P,\operatorname{Cone}(\mathsf P\to P))\);
a finite complex of projectives takes an acyclic complex to an acyclic
Hom complex, by induction on its finite length.
The resulting Hom differential is zero since \(pP=0\).
Its two terms are \(A\) in degrees zero and one. The composite from
\(E\) is an isomorphism onto these terms. This proves the claimed
quasi-isomorphism of algebras, including the Yoneda multiplication.
\end{proof}

Thus replacing the derived section by the underived algebra \(A\)
would delete an actual extension class. Adding a cotangent coordinate
and then taking a section is not an inverse operation on derived
endomorphisms.

\section{A central surface in a noncommutative algebra}

Introduce a parameter \(\lambda\), independent of the fixed Yangian
parameter \(\hbar\), and define
\begin{equation}\label{eq:cs-deformation}
 B_\lambda=
 \mathbb C[\lambda]\langle z,p,g\rangle/
 (g^2-1,\ gz+zg,\ gp+pg,\ pz-zp-\lambda g).
\end{equation}
This means one algebra over \(\mathbb C[\lambda]\); evaluation at a
complex number gives its corresponding fibre.

\begin{lemma}\label{lem:cs-pbw}
The words \(z^ip^jg^e\), \(i,j\geq0\), \(e=0,1\), form a free
\(\mathbb C[\lambda]\)-basis of \(B_\lambda\).
\end{lemma}
\begin{proof}
Rewrite \(pz\) as \(zp+\lambda g\), \(gz\) as \(-zg\),
\(gp\) as \(-pg\), and \(gg\) as \(1\).
Order a word first by its length and then by its number of reversed
pairs for the order \(z<p<g\). Each replacement decreases this order
in its non-leading terms and decreases the reversed-pair count in
its equal-length term. Thus reduction terminates.
The overlapping reductions are \(gpz,ggz,ggp,ggg\).
Both reductions of \(gpz\) give \(zpg+\lambda\); the other three
agree by cancellation of two signs or two copies of \(g\).
Reductions on disjoint subwords commute. Induction on the terminating
order now proves uniqueness of the reduced expression: any first
two reductions can be joined, and their smaller continuations agree.
The vector space of reduced words, with multiplication by reduction,
is consequently associative and realizes the presentation. This
proves independence as well as spanning.
\end{proof}

Put
\begin{equation}\label{eq:cs-central-coordinates}
 X=z^2,\qquad Y=p^2,\qquad W=zp+\lambda g/2.
\end{equation}

\begin{theorem}\label{thm:cs-centre}
The centre of \(B_\lambda\) and the corner \(e_+B_\lambda e_+\) are
both the algebra
\[
 C_\lambda=
 \mathbb C[\lambda,X,Y,W]/(XY-W^2+\lambda^2/4).
\]
The map to the corner sends a central element to its product with
\(e_+\).
\end{theorem}
\begin{proof}
The relation \([p,z^2]=\lambda(gz+zg)=0\) proves centrality of
\(X\); the same argument applies to \(Y\). Direct commutation with
\(z,p,g\) proves centrality of \(W\). Expanding its square gives
\(W^2=z^2p^2+\lambda^2/4\).

A word commuting with \(g\) has even total \(z,p\) degree.
If both exponents are even it is a monomial in \(X,Y\), possibly
times \(g\). If both are odd, replace \(zp\) by \(W-\lambda g/2\).
Thus the even subspace is \(C_\lambda\oplus C_\lambda g\).
There are no hidden relations: reduce a polynomial in \(W\) to
degree at most one by the monic equation for \(W^2\).
The monomials \(X^aY^b\) and \(X^aY^bW\), with or without \(g\),
have distinct highest normal words in Lemma~\ref{lem:cs-pbw}.
They are independent.

If \(F+Hg\) is central, its commutator with \(z\) is \(2Hzg\).
Left multiplication by \(z\) is injective in the normal basis.
Since \(H\) is central, this forces \(H=0\). This proves the centre.
In the corner, odd-total-degree words vanish and \(g\) acts as \(1\).
The same independence proves the stated corner isomorphism.
\end{proof}

In terms of the four arrows, \eqref{eq:cs-deformation} reads
\begin{equation}\label{eq:cs-deformed-arrows}
 cb-ad=\lambda e_+,\qquad da-bc=-\lambda e_-.
\end{equation}
After inverting \(\lambda\), the second equation gives
\begin{equation}\label{eq:cs-full}
 1=e_++\lambda^{-1}be_+c-\lambda^{-1}de_+a.
\end{equation}
This is a full-idempotent identity, so the multiplication maps from
Chapter~\ref{chap:fc-fusion} give a Morita equivalence between
\(B_\lambda[\lambda^{-1}]\) and \(C_\lambda[\lambda^{-1}]\).
It applies to modules and complexes through explicit tensor functors.
At \(\lambda=0\), the one-dimensional module with \(z=p=0\) and
\(g=-1\) is nonzero and is killed by \(e_+\). A full idempotent
cannot kill a nonzero module, since applying
\(1=\sum x_i e_+ y_i\) would then kill every vector.
The equivalence therefore fails at that parameter.

\section{The bundle which retains both fusion modules}

Fix a scalar \(\lambda\ne0\). The affine surface
\(\mathcal S_\lambda\) has coordinate ring \(C_\lambda\) at this
fixed parameter. It is smooth: a singular point would satisfy
\(X=Y=W=0\) by the three partial derivatives of its equation, and
that point does not satisfy the equation when \(\lambda\ne0\).

A line bundle on this surface may be specified by a module which
becomes free of rank one on an open cover. Consider the ideal
\begin{equation}\label{eq:cs-line}
 I=(X,W-\lambda/2)\subset C_\lambda.
\end{equation}
On the open set \(U_+\) where \(W+\lambda/2\ne0\), the surface
equation gives \(I=(X)\). On \(U_-\), where \(W-\lambda/2\ne0\),
it gives \(I=C_\lambda\). These sets cover since \(\lambda\ne0\).
Multiplication by \(X\) is injective, so \(X\) is a free generator
on the first chart. Thus \(I\) is an invertible module.

Left multiplication by \(a\) identifies \(e_-B_\lambda e_+\)
with \(I\). Indeed this module is generated over the centre by
\(b,d\), by reducing odd normal words. Their images are
\(ab=X\) and \(ad=W-\lambda/2\).
If \(a\ell=0\), then \(ba\ell=X\ell=0\). Injectivity of
multiplication by \(z^2\) in the normal basis gives \(\ell=0\).
This proves the module identification.

Consequently \(B_\lambda e_+\) is \(C_\lambda\oplus I\).
Its left action is
\begin{equation}\label{eq:cs-bundle-action}
 \begin{split}
 z(c,i)&=(i,Xc),\\
 p(c,i)&=((W+\lambda/2)i/X,\ (W-\lambda/2)c),\\
 g(c,i)&=(c,-i).
 \end{split}
\end{equation}
The apparently rational entry is a regular map \(I\to C_\lambda\):
it sends the two generators \(X,W-\lambda/2\) to
\(W+\lambda/2,Y\). In the local bases \((1,X)\) on \(U_+\)
and \((1,1)\) on \(U_-\), these actions are
\[
 z_+=\begin{pmatrix}0&X\\1&0\end{pmatrix},\qquad
 p_+=\begin{pmatrix}0&W+\lambda/2\\Y/(W+\lambda/2)&0\end{pmatrix},
\]
\[
 z_-=\begin{pmatrix}0&1\\X&0\end{pmatrix},\qquad
 p_-=\begin{pmatrix}0&Y/(W-\lambda/2)\\W-\lambda/2&0\end{pmatrix}.
\]
On the overlap \(X\) is invertible, since
\(XY=(W-\lambda/2)(W+\lambda/2)\).
The transition \(T=\operatorname{diag}(1,X)\) gives
\(z_-=Tz_+T^{-1}\) and \(p_-=Tp_+T^{-1}\).
The matrices satisfy \(z^2=X\), \(p^2=Y\), and
\([p,z]=\lambda\operatorname{diag}(1,-1)\).
The full-idempotent identity proves that this action identifies
\[
 B_\lambda\simeq\operatorname{End}_{C_\lambda}(C_\lambda\oplus I).
\]

The divisor \(X=0\) has two disjoint components
\[
 D_+=\{X=0,\ W=\lambda/2\},\qquad
 D_-=\{X=0,\ W=-\lambda/2\}.
\]
On \(D_+\), the \(z\)-arrows have \(a=X,b=1\); on \(D_-\), they
have \(a=1,b=X\). These are the opposite fusion modules
of \eqref{eq:fc-probes}. The smoothing has placed both on one
surface while preserving their difference.

\begin{theorem}\label{thm:cs-yangian-bundle}
Retain the local fusion algebra \(\Lambda_d\), its full corner
\(A_t\), and its selected idempotent \(e\) from
Chapter~\ref{chap:fc-fusion}. On the open surface
\(X+2\hbar\ne0\), with coefficient map \(t\mapsto X\), the module
\[
 \mathcal V_\lambda
 =\Lambda_de\otimes_{A_t}B_\lambda e_+
 \simeq \mathcal O_{\mathcal S_\lambda}^{\oplus p}
          \oplus I^{\oplus q}
\]
is a vector bundle of rank \(d^2\) with the original tensor
Yangian action. Its two restrictions at \(X=0\) retain the two
fusion limits.
\end{theorem}
\begin{proof}
The map \(A_t\to B_\lambda\) sends its two arrows to the \(z\)-arrows,
and sends \(t\) to \(X\). Their products are \(X e_\pm\), so it
is an algebra map. Localizing at \(X+2\hbar\) is exactly the
localization used in the fusion coefficient ring.
The Yangian acts on \(\Lambda_de\) on the left, commuting with its
right \(A_t\)-action, and hence acts on the tensor product.
As a right \(A_t\)-module, \(\Lambda_de\) is a sum of \(p\) copies
of \(e_0A_t\) and \(q\) copies of \(e_1A_t\).
Tensoring gives the stated copies of \(C_\lambda\) and \(I\).
Their local bases prove the rank and the two specializations.
\end{proof}

This action does not assert that the Yangian alone generates all
endomorphisms on the collision components. There its image preserves
the appropriate flag. The added \(p\)-operators generate the larger
surface algebra. The two parameters also have different roles:
\(\hbar\) specifies the evaluation representation, while \(\lambda\)
deforms the cotangent algebra.

\section{Cyclic ramification and the exact smoothness condition}

Let \(r\geq2\) and let \(\zeta\) be a primitive \(r\)-th root of
unity. The cyclic group acts on \(\mathbb C[z]\) by \(g(z)=\zeta z\).
Its character idempotents are
\[
 e_i=\frac1r\sum_{j=0}^{r-1}\zeta^{-ij}g^j,\qquad i\in\mathbb Z/r.
\]
They are orthogonal and sum to \(1\), by summing a finite geometric
series. Multiplication by \(z\) sends vertex \(i\) to \(i+1\).
Over \(\mathbb C[X]\), \(X=z^r\), the skew group algebra
\(\mathbb C[z]\rtimes\mu_r\) is free of rank \(r^2\).
Each full cyclic path equals \(X\). These equations are the cyclic
version of the quadratic cover.

Choose complex numbers \(\kappa_0,\ldots,\kappa_{r-1}\), use cyclic
indices, and put \(\lambda_i=\kappa_i-\kappa_{i+1}\).
Define \(B_{\boldsymbol\kappa}\) by the idempotents and two arrows
at each vertex:
\begin{equation}\label{eq:cs-cyclic-algebra}
 e_{i+1}z=ze_i,\qquad e_{i-1}p=pe_i,\qquad
 pz-zp=\sum_i\lambda_i e_i.
\end{equation}
When all \(\kappa_i\) coincide this is
\(\mathbb C[z,p]\rtimes\mu_r\), with action
\((z,p)\mapsto(\zeta z,\zeta^{-1}p)\).
The integer \(r\) counts characters of this cover. It does not
count Yangian evaluation factors.

\begin{lemma}\label{lem:cs-cyclic-normal}
The elements \(z^ap^be_i\), \(a,b\geq0\), \(0\leq i<r\), are a
basis of \(B_{\boldsymbol\kappa}\).
\end{lemma}
\begin{proof}
On \(A=\mathbb C[z]\rtimes\mu_r\), set
\(\sigma(z)=z\), \(\sigma(g)=\zeta g\).
This is an automorphism. Define a \(\sigma\)-derivation by
\(\delta(z)=\sum_i\lambda_i e_i\), \(\delta(g)=0\), and
\(\delta(ab)=\sigma(a)\delta(b)+\delta(a)b\).
It respects the relations: the derivatives of \(gz\) and
\(\zeta zg\) are both
\(\zeta g\sum_i\lambda_i e_i\), and the derivative of \(g^r-1\)
is zero. Thus it is well-defined on \(A\).
Adjoin \(p\) by the rule \(pa=\sigma(a)p+\delta(a)\).
This rule determines multiplication uniquely on the free left
\(A\)-module with basis \(1,p,p^2,\ldots\).
The \(\sigma\)-derivation identity verifies
\(p(ab)=(pa)b\); induction on the power of \(p\) proves
associativity. It therefore constructs the presentation with no
further linear relations. Moving group elements to the right and
using the character basis gives the claimed normal basis.
\end{proof}

Define
\[
 X=z^r,\qquad Y=p^r,\qquad
 W=zp+\sum_i\kappa_i e_i,\qquad
 P(W)=\prod_{i=0}^{r-1}(W-\kappa_i).
\]

\begin{theorem}\label{thm:cs-cyclic-centre}
These three elements are central, and
\[
 Z(B_{\boldsymbol\kappa})
 \simeq e_0B_{\boldsymbol\kappa}e_0
 \simeq C_{\boldsymbol\kappa}
 :=\mathbb C[X,Y,W]/(XY-P(W)).
\]
The map to the corner is multiplication by \(e_0\).
\end{theorem}
\begin{proof}
At vertex \(i\) the defining equations give
\[
 zp\,e_i=(W-\kappa_i)e_i,\qquad
 pz\,e_i=(W-\kappa_{i+1})e_i.
\]
These show that \(W\) commutes with \(z,p\).
Moving \(p\) past \(z^r\) sums the differences
\(\kappa_i-\kappa_{i+1}\) around a full cycle, giving zero.
The same argument proves centrality of \(p^r\).
Iterating the displayed products gives \(XY=P(W)\).

A corner word has \(a-b\) divisible by \(r\). Removing powers
of \(X,Y\) leaves a product \(z^np^n\), and the displayed equations
express this as a polynomial in \(W\). Thus the proposed ring
surjects onto the corner. Since \(P\) is monic of degree \(r\),
the quotient has basis \(X^aY^bW^c\), \(a,b\geq0\), \(0\leq c<r\),
over \(\mathbb C\): divide by the monic polynomial in \(W\).
Their images have distinct highest normal words
\(z^{ra+c}p^{rb+c}e_0\), so are independent.

The same argument at every vertex shows that the centralizer of
the character idempotents is \(\bigoplus_i C_{\boldsymbol\kappa}e_i\).
Commutation with \(z\) makes the neighboring coefficients equal,
since multiplication by \(z\) is injective on the normal basis.
All coefficients are therefore equal. This proves the centre claim.
\end{proof}

\begin{theorem}\label{thm:cs-full-iff}
The following are equivalent:
\begin{enumerate}
\item the numbers \(\kappa_0,\ldots,\kappa_{r-1}\) are pairwise distinct;
\item the surface \(XY=P(W)\) is smooth;
\item \(e_0\) is a full idempotent of \(B_{\boldsymbol\kappa}\).
\end{enumerate}
When they hold, every vertex idempotent is full.
When two roots coincide, an explicit nonzero finite-dimensional
module is annihilated by \(e_0\).
\end{theorem}
\begin{proof}
A singular point of the surface must satisfy
\(X=Y=0\), \(P(W)=P'(W)=0\).
Such a point exists precisely when \(P\) has a repeated root.
This proves the first equivalence.

Suppose the roots are distinct. For \(1\leq j<r\), put
\[
 A_j(W)=\prod_{i=1}^j(W-\kappa_i),\qquad
 B_j(W)=\prod_{i=j+1}^r(W-\kappa_i),\qquad \kappa_r=\kappa_0.
\]
These polynomials have no common root. The Euclidean algorithm
therefore supplies polynomials \(\alpha_j,\beta_j\) with
\(\alpha_jA_j+\beta_jB_j=1\).
The path products give
\[
 (z^je_0)(p^je_j)=A_j(W)e_j,\qquad
 (p^{r-j}e_0)(z^{r-j}e_j)=B_j(W)e_j.
\]
Consequently
\begin{equation}\label{eq:cs-bezout}
 e_j=\alpha_j(W)(z^je_0)(p^je_j)
       +\beta_j(W)(p^{r-j}e_0)(z^{r-j}e_j).
\end{equation}
Together with \(e_0\), their sum expresses \(1\) in
\(B_{\boldsymbol\kappa}e_0B_{\boldsymbol\kappa}\).
Cyclic relabeling proves fullness of every vertex.

Conversely, suppose two roots coincide. Order them in the linear
list \(\kappa_1,\ldots,\kappa_r\), where \(\kappa_r=\kappa_0\).
There are \(1\leq a\leq b<r\) with \(\kappa_a=\kappa_{b+1}\).
Put a one-dimensional vector space at each vertex \(a,\ldots,b\)
and zero at all other vertices. For \(a\leq i<b\), let the
\(z\)-arrow \(i\to i+1\) be \(1\), and the reverse \(p\)-arrow
be \(\kappa_a-\kappa_{i+1}\).
All remaining arrows are zero.
At an interior vertex their two products differ by
\[
 (\kappa_a-\kappa_{i+1})-(\kappa_a-\kappa_i)
 =\kappa_i-\kappa_{i+1}.
 \]
At the first vertex the second term is zero, and the same equation
holds. At the last vertex the first term is zero; the equality
\(\kappa_a=\kappa_{b+1}\) gives the required equation there as well.
For \(a=b\) both arrows are zero and
\(\lambda_a=\kappa_a-\kappa_{a+1}=0\).
Thus all cases define a representation of
\eqref{eq:cs-cyclic-algebra}. It is nonzero and has no space at
vertex zero, so \(e_0\) kills it.
The full-idempotent identity would force that module to vanish.
This proves necessity and the final assertion.
\end{proof}

For example, at \(r=3\), if \(\kappa_1=\kappa_2\), the missing
module can be supported at vertex \(1\) alone. If instead
\(\kappa_1=\kappa_0\), it is supported at vertices \(1,2\), with
forward arrow \(1\) and reverse arrow \(\kappa_1-\kappa_2\).
The repeated-root condition is detected by actual representations,
including cases where the chosen root \(\kappa_0\) itself remains
distinct from the two coincident roots.

\section{Line ideals and every character module}

Assume the roots are distinct, and put
\[
 I_0=C_{\boldsymbol\kappa},\qquad
 I_j=(X,B_j(W))\quad(1\leq j<r),
\]
with \(B_j\) as in the proof of Theorem~\ref{thm:cs-full-iff}.
Multiplication on the left by \(z^{r-j}e_j\) identifies
\(e_jB_{\boldsymbol\kappa}e_0\) with \(I_j\).
Indeed that module is generated over the centre by the two paths
\(z^je_0,p^{r-j}e_0\). Their images are \(X,B_j(W)\).
Injectivity follows from injectivity of multiplication by \(z\)
in the normal basis.

On \(A_j(W)\ne0\), the equation \(XY=A_jB_j\) gives \(I_j=(X)\).
On \(B_j(W)\ne0\), it gives \(I_j=C_{\boldsymbol\kappa}\).
The Bezout identity shows that these open sets cover.
Thus each \(I_j\) is a line bundle and
\begin{equation}\label{eq:cs-all-bundle}
 B_{\boldsymbol\kappa}\simeq
 \operatorname{End}_{C_{\boldsymbol\kappa}}
       (I_0\oplus I_1\oplus\cdots\oplus I_{r-1}).
\end{equation}
This is the left action on \(B_{\boldsymbol\kappa}e_0\).
The full-idempotent multiplication identities prove that it is
onto every endomorphism, as well as injective.

The local matrices make the characters visible. The open sets
\[
 U_k=\{P_k(W)\ne0\},\qquad
 P_k(W)=\prod_{i\ne k}(W-\kappa_i),\qquad 0\leq k<r,
\]
cover the surface: away from all roots every \(P_k\) is nonzero,
and at a root exactly its own \(P_k\) is nonzero.
On \(U_k\), choose the generator \(X\) of \(I_j\) if \(\kappa_k\)
is a root of \(B_j\), and choose the generator \(1\) otherwise.
Write \(v_i\) for these local basis vectors, including \(v_0=1\).
Then
\begin{equation}\label{eq:cs-all-matrices}
 zv_i=c_i v_{i+1},\qquad
 pv_{i+1}=\frac{W-\kappa_{i+1}}{c_i}v_i,\qquad
 c_i=\begin{cases}X&i+1=k\pmod r,\\1&\text{otherwise.}\end{cases}
\end{equation}
The sole apparent denominator \(X\) is harmless because
\((W-\kappa_k)/X=Y/P_k(W)\) is regular on this chart.
To derive the \(z\)-coefficients, use the embeddings
\(e_jBe_0\to I_j\): for \(1\leq j<r-1\), the image of \(z\ell\)
in \(I_{j+1}\) is the same element as the image of \(\ell\)
in \(I_j\). The first step multiplies by \(X\), and the last
step includes \(I_{r-1}\) in \(I_0\).
Taking the ratios of the chosen generators gives exactly \(c_i\).
The equation \(pz\,e_i=(W-\kappa_{i+1})e_i\) then gives \(p\).
These operators satisfy \(z^r=X\), \(p^r=Y\), and
\([p,z]v_i=(\kappa_i-\kappa_{i+1})v_i\).
The transition maps are ratios of local generators of the same
ideals, so the matrices glue.

At \(X=0\) there are \(r\) disjoint components \(W=\kappa_k\).
On the \(k\)-th component exactly the arrow \(k-1\to k\)
vanishes, while the other \(z\)-arrows are \(1\).
These are the \(r\) distinct character modules of the original
cyclic cover: in each equivariant free line, multiplication by
\(z\) moves through the character eigenspaces, and one full
cycle multiplies by \(X\). Changing the character changes the
position of that last arrow. Thus every character has a specified
geometric continuation.

The two-form
\begin{equation}\label{eq:cs-volume}
 \omega=\frac1r\,\frac{dX\wedge dW}{X}
\end{equation}
on \(X\ne0\) extends across the smooth surface. On \(Y\ne0\)
it is \(-dY\wedge dW/(rY)\), and on \(P'(W)\ne0\) it is
\(dX\wedge dY/(rP'(W))\). Differentiating \(XY=P(W)\) verifies
their agreement. These charts cover the smooth surface, and each
expression is nonvanishing in its local coordinates.
At the undeformed cover \(X=z^r,W=zp+\kappa_0\), its pullback
is \(dz\wedge dp\). The factor \(1/r\) fixes that normalization.

\section{A self-dual diagonal at every deformation parameter}

The smooth surface describes the module category when the roots are
distinct. The algebra itself has a two-dimensional duality even when
some roots coincide. To prove this, we construct its bimodule resolution
and its dual, rather than infer duality from a volume form on its centre.

Write \(B=B_{\boldsymbol\kappa}\). Denote the arrow \(i\to i+1\)
by \(z_i\), its reverse by \(p_i\), and set
\(z_i^*=p_i\), \(p_i^*=z_i\).
Put \(\varepsilon(p_i)=1\), \(\varepsilon(z_i)=-1\).
If \(h(a)\) and \(t(a)\) are the head and tail of an arrow, the
relations are
\[
 \sum_{h(a)=i}\varepsilon(a)aa^*=\lambda_i e_i.
\]
Introduce the projective \(B\)-bimodules
\[
 P_0=\bigoplus_i Be_i\otimes_{\mathbb C}e_iB,\qquad
 P_1=\bigoplus_a Be_{h(a)}\otimes_{\mathbb C}e_{t(a)}B,\qquad
 P_2=P_0.
\]
Let \(\eta_i\) and \(\eta_a\) be the indicated tensor generators.
Each summand is an idempotent summand of the free bimodule
\(B\otimes_{\mathbb C}B\). Define
\begin{align}\label{eq:cs-diagonal}
 d_1\eta_a&=a\eta_{t(a)}-\eta_{h(a)}a,\\
 d_2\eta_i&=\sum_{h(a)=i}\varepsilon(a)
                      (\eta_a a^*+a\eta_{a^*}),\qquad
 \mu(\eta_i)=e_i.\notag
\end{align}
The augmentation \(\mu:P_0\to B\) is multiplication.
The square \(d_1d_2\eta_i\) is
\(\lambda_i e_i\otimes e_i-e_i\otimes\lambda_i e_i=0\).
Thus these maps form an augmented complex.

\begin{proposition}\label{prop:cs-self-dual}
For all values of \(\boldsymbol\kappa\),
\[
 0\longrightarrow P_2\xrightarrow{d_2}P_1
 \xrightarrow{d_1}P_0\xrightarrow{\mu}B\longrightarrow0
\]
is exact. Moreover
\[
 \operatorname{RHom}_{B-B}(B,B\otimes_{\mathbb C}B)\simeq B[-2]
\]
as bimodule complexes, with the inner bimodule action on Hom.
In particular \(B\) is a smooth two-Calabi--Yau algebra in the
bimodule sense.
\end{proposition}
\begin{proof}
Filter by path length, giving the displayed generators of \(P_n\)
weight \(n\). The associated graded algebra is
\(\mathbb C[z,p]\rtimes\mu_r\) by
Lemma~\ref{lem:cs-cyclic-normal}. The associated graded diagonal
complex is the two-variable Koszul diagonal with the character
decomposition retained.
Here is its exactness explicitly. Put \(S=\mathbb C[z,p]\).
The map
\[
 (s g^i)\otimes_{\mathbb C[\mu_r]}(t g^j)
 \longmapsto s\otimes g^i(t)\otimes g^{i+j}
\]
identifies the degree-zero term with
\((S\otimes S)\otimes\mathbb C[\mu_r]\).
In degree one include the two symbols \(dz,dp\), transforming
by \(\zeta,\zeta^{-1}\); in degree two include
\(dz\wedge dp\), which is invariant.
The same formula transforms these symbols along with \(t\).
It identifies \eqref{eq:cs-diagonal} with the Koszul differential
on \(z_L-z_R,p_L-p_R\), tensor the group algebra.
These two differences are independent polynomial variables.
If \(Za+Pb=0\), reduction modulo \(Z\) gives \(b=Zc\),
and then \(a=-Pc\). Multiplication by the pair \((P,-Z)\)
is injective. The cokernel of the row \((Z,P)\) is the
polynomial diagonal. This proves exactness of the associated
graded complex, with a harmless choice of sign for its top generator.
For the filtered complex, take a cycle of greatest finite weight.
Lift a primitive for its leading term, subtract its differential,
and repeat at smaller weight. This proves exactness; injectivity
at the top follows from injectivity of its associated graded map.

It remains to identify the dual complex. For a bimodule map with
value \(\sum u\otimes v\), the inner action is
\((b\phi c)(m)=\sum uc\otimes bv\).
If \(\phi(\eta_i)=\sum u_i\otimes v_i\), define
\[
 \alpha(\phi)=\sum_i v_i\eta_i u_i.
\]
For \(\psi(\eta_a)=\sum u_a\otimes v_a\), define
\[
 \beta(\psi)=\sum_a\varepsilon(a)v_a\eta_{a^*}u_a.
\]
The sums also run over the tensor terms in each value.
These are bimodule isomorphisms: their inverses reverse the two
outer factors, and \(\beta^{-1}\) also reverses the arrow and its
specified sign.
The Hom differential is \(-d_1^*\) in degree zero and \(d_2^*\)
in degree one. Direct substitution gives the two identities
\[
 \beta(-d_1^*\phi)=d_2\alpha(\phi),\qquad
 \alpha(d_2^*\psi)=d_1\beta(\psi).
\]
For the first, the head terms are
\(\varepsilon(a)v_{h(a)}a\eta_{a^*}u_{h(a)}\);
the tail terms become the other summands of \(d_2\alpha\)
after replacing \(a\) by \(a^*\) and using
\(\varepsilon(a^*)=-\varepsilon(a)\).
For the second, expand the two terms
\(\eta_a a^*\) and \(a\eta_{a^*}\) in \(d_2\eta_i\).
Reversing their outer factors gives respectively
\(\varepsilon(a)v_a a^*\eta_{h(a)}u_a\) and
\(-\varepsilon(a)v_a\eta_{t(a)}a^*u_a\), exactly
\(d_1\beta(\psi)\).
Thus the dual resolution is the original three-term complex
placed in cochain degrees \(0,1,2\).
Its only cohomology is \(B\) in degree two, proving the formula.
The finite projective resolution proves homological smoothness.
These two assertions are the stated bimodule definition of the
two-Calabi--Yau property.
\end{proof}

No finite-dimensionality or properness is asserted here.
At repeated roots the centre is singular, but the algebra still
has the displayed finite self-dual resolution. Its additional
character modules are therefore part of the algebraic geometry
that the singular centre alone does not retain.

\section{What changes at the zero section}

For the quadratic deformation at \(\lambda\ne0\), no quotient
can retain \(p=0\): the relation \(pz-zp=\lambda g\) would give
\(g=0\), contradicting \(g^2=1\) in a nonzero unital quotient.
The zero-section module from
Proposition~\ref{prop:cs-normal-extension} therefore cannot be
held fixed through this deformation.
The projective module \(B_\lambda e_+\), on the other hand, is
defined throughout the family. It becomes a full generator
precisely on the nonzero-parameter locus. At zero it misses the
sign module. These are two different changes: a supported module
ceases to exist, while a projective module changes its generating
property.

The same distinction is now exact in every cyclic order.
Theorem~\ref{thm:cs-full-iff} identifies the singular parameter
locus and supplies the modules lost by a one-vertex description.
On the complement, \eqref{eq:cs-all-bundle} gives the full
endomorphism algebra of an explicit bundle. This is an equivalence
in that parameter region, not an equivalence between the different
fibres of a deformation.

Finally, the independent cyclic cover must not be substituted for
an overlapping Yangian collision. On three tensor factors,
\[
 P_{12}P_{23}(e_0\otimes e_0\otimes e_1)
   =e_1\otimes e_0\otimes e_0,\qquad
 P_{23}P_{12}(e_0\otimes e_0\otimes e_1)
   =e_0\otimes e_1\otimes e_0.
\]
These exchanges do not commute. Constructing a common integral
lattice at overlapping resonances requires their simultaneous
relations. Neither the cyclic cover nor a tensor product of
independent covers supplies that additional construction.
